\documentclass{tutodoc}

\usepackage{../preamble.cfg}

\usepackage{../main/core-ctxt-dsl}
\usepackage{../main/main}


% TESTING LOCAL IMPLEMENTATION %

\usepackage{data-1D}


\begin{document}

\section{Tableaux d'images de fonctions à une seule variable}

La saisie de tableaux, automatiquement en mode mathématique, d'images de fonctions à une seule variable nécessitent de fournir les valeurs pivots, c'est-à-dire les nombres dont on veut indiquer les images. Ceci peut se faire de deux façons.
%
\begin{enumerate}
	\item \tdoclatexin{xvals = x1 , x2 , ... , xn} indiquent les valeurs pour la variable $x$, qui est celle par défaut.

	\item \tdoclatexin{xvals = mavar : x1 , x2 , ... , xn} permet de renommer la variable en $mavar$.
\end{enumerate}


Pour les images, il existe deux modes de saisie, mais avec une restriction forte, car une seule série d'images peut utiliser le nom par défaut.
%
\begin{enumerate}
	\item \tdoclatexin{xvals = x1 , x2 , ... , xn} indiquent les images pour la fonction $f(x)$, qui est celle par défaut.

	\item \tdoclatexin{imgs = mafonc : i1 , i2 , ... , in} indique les images pour la fonction $mafonc$.
\end{enumerate}

En résumé, tout se passe via \tdoclatexin{xvals = ...} pour les valeurs pivots, et \tdoclatexin{imgs = ...} pour les images.

% --------------- %


\begin{tdocexa}[Une seule fonction avec les noms par défaut]
    \leavevmode

    \tdoclatexinput{examples/data-1D/one-func-no-label.tex}
\end{tdocexa}


\begin{tdocnote}
	Retenir que tout se saisie en mode mathématique.
\end{tdocnote}


\begin{tdocwarn}
	\begin{enumerate}[wide]
		\item L'utilisation de \tdoclatexin{xvals} doit être faite une fois, et une seule, au tout début du contenu.

		\item L'emploi du nom par défaut pour une fonction n'est possible que pour une seule série d'images.
	\end{enumerate}
\end{tdocwarn}


% --------------- %


\begin{tdocexa}[Deux fonctions et une variable, le tout en mode \tdocquote{maison}]
    \leavevmode

    \tdoclatexinput{examples/data-1D/two-func-label.tex}
\end{tdocexa}


% --------------- %


\begin{tdocexa}[Commentaires à la sauce \LaTeX]
    \leavevmode

    \tdoclatexinput{examples/data-1D/comments.tex}
\end{tdocexa}


% --------------- %

\begin{tdoctip}[Nombres décimaux en version \tdocquote{locale} et \tdocquote{grandes} fractions]
%    \leavevmode
%
    Via les macros \tdocmacro{dfrac} et \tdocmacro{num}%
    \footnote{
    	Cette macro ajoute de fins espaces mettant en valeur les groupes de chiffres, tout en gérant l'absence d'espaces autour du séparateur décimal.
    }
    venant des excellents packages \tdocpack{amsmath} et \tdocpack{siuntix} respectivement, il est facile de rédiger des nombres décimaux, et d'obtenir de \tdocquote{grandes} fractions,%
    \footnote{
    	Bien que typographiquement laides, les grandes fractions peuvent être très utiles dans un contexte pédagogique de niveau pré-universitaire.
    }
    comme le montre l'exemple suivant.

    \tdoclatexinput{examples/data-1D/num-n-dfrac.tex}
\end{tdoctip}

\end{document}
